\subsection{The exact simultaneous geometric and tangential image}

\begin{result}{E.127}{Proved}{The exact joint image of the point $5/16$ and all actual tangential actions}
Use the integral unit tangent $\overrightarrow{01}$, and write
\[
 \rho_\ell:G_{\mathbb Q}\longrightarrow\operatorname{Aut}(F_\ell),
 \qquad
 \rho_{\rm all}=\prod_\ell\rho_\ell,
 \qquad D=\rho_{\rm all}(G_{\mathbb Q}),\quad K=\ker\rho_{\rm all}.
\]
Let $t=5/16$, and let
$z_t:G_{\mathbb Q}\to\prod_\ell F_\ell$
be the point cocycle obtained from any one geometric path.
Then
\[
 \prod_{\ell\ \mathrm{odd}}F_\ell
        \ \subseteq\ z_t(K)\ \subseteq\ \prod_\ell F_\ell,
                                                        \tag{E.127.1}
\]
where the left-hand group has identity two-primary coordinate.
Consequently the actual joint homomorphism
\[
 G_{\mathbb Q}\longrightarrow
       \left(\prod_{\ell\ \mathrm{odd}}F_\ell\right)\rtimes D,
 \qquad g\longmapsto\bigl(z_{t,\mathrm{odd}}(g),\rho_{\rm all}(g)\bigr)
                                                        \tag{E.127.2}
\]
is surjective.  More precisely, if $H_2=z_t(K)\cap F_2$, then
\[
 H_2=\Phi(F_2),\qquad
 z_t(K)=\Phi(F_2)\times\prod_{\ell\ \mathrm{odd}}F_\ell.
                                                        \tag{E.127.3}
\]
Here $\Phi(F_2)=\overline{F_2^2[F_2,F_2]}$ is its
Frattini subgroup.  The two quadratic Kummer characters of
$5$ and $11$ descend to a continuous homomorphism
$\nu:D\to\mathbb F_2^2$.  The image of the full joint map
$g\mapsto(z_t(g),\rho_{\rm all}(g))$ is exactly
\[
 \left\{(w,d)\in\left(\prod_\ell F_\ell\right)\rtimes D:
          w_2\bmod\Phi(F_2)=\nu(d)\right\}.
                                                        \tag{E.127.4}
\]
It has index four in this semidirect product.
\end{result}
\begin{proof}
The actual tangential action $\rho_\ell$ is unramified at
every finite rational prime other than $\ell$.  This is
Theorem~2.1 of Wojtkowiak, with its Definition~2.0: the
marked pair $(U,\overrightarrow{01})$ has good reduction
at every finite prime \cite[pp.198--200]{AuditWojtkowiak2007}.
This assertion is about automorphisms with the specified
basepoint, so inertia acts as the identity, not merely as
an inner automorphism.

We also need a group-theoretic fact with its exact scope.
For a finitely generated pro-$q$ group $P$, the kernel of
\[
       \operatorname{Aut}(P)\longrightarrow
                         \operatorname{Aut}(P/\Phi(P))
\]
is pro-$q$.  To see this, pass to characteristic finite
$q$-group quotients.  For a finite $q$-group, the kernel
acts trivially on the associated graded of the lower
exponent-$q$ central series, since that graded group is
generated by the first quotient under commutators and
$q$th powers.  Filtering automorphisms by their first
possible nontrivial displacement gives successive
subgroups of additive homomorphism groups between finite
$\mathbb F_q$-vector spaces.  Each finite kernel is
therefore a $q$-group, and the inverse limit is pro-$q$.
In particular, the kernel of the action on the full
abelianization of $F_q$ is pro-$q$.

Choose one complex conjugation $c\in G_{\mathbb Q}$.
For each $q\in\{2,5,11\}$, choose a tame inertia element
$\tau_q\in I_q$ with a topological generator as its
prime-to-$q$ tame image.  Put
\[
                         h_q=[c,\tau_q].
\]
If $r\ne q$, then $\rho_r(\tau_q)=1$ and hence
$\rho_r(h_q)=1$.  The action of $G_{\mathbb Q}$ on
$F_q^{\rm ab}=\mathbb Z_q(1)^2$ is scalar cyclotomic.
Its commutator is therefore trivial on that
abelianization.  It follows that $\rho_q(h_q)$ belongs
to the pro-$q$ kernel just described.  Thus the entire
image $\rho_{\rm all}(h_q)$ is supported at $q$ and is
contained in a pro-$q$ group.

For a prime $\ell\ne q$, let $e_\ell\in\widehat{\mathbb Z}$
have $\ell$-component one and every other component zero,
and set $k_{q,\ell}=h_q^{e_\ell}$.  Profinite powers are
taken in the procyclic closure of $h_q$.  Its
$e_\ell$-part is pro-$\ell$; its image in the stated
pro-$q$ arithmetic group is trivial.  Consequently
\[
              k_{q,\ell}\in K,\qquad
              z_t(k_{q,\ell})\in F_\ell.               \tag{E.127.5}
\]
The second assertion follows because $z_t|_K$ is a
homomorphism and the image of a pro-$\ell$ group in a
pronilpotent group lies in its $\ell$-factor.

We compute the images modulo $\Phi(F_\ell)$.  Let
$\kappa_\ell=(\kappa_\ell(t),\kappa_\ell(1-t))$ be
the abelian Kummer cocycle of the point.  For any
scalar-character cocycle $a$ one has
\[
 a([g,h])=(1-\chi(h))a(g)+(\chi(g)-1)a(h).
\]
Here $\chi_\ell(\tau_q)=1$ for $q\ne\ell$, and
$\chi_\ell(c)=-1$.  The tame Kummer formula therefore
gives
\[
 \kappa_\ell(h_q)=-2u_{q,\ell}\,
             \bigl(v_q(t),v_q(1-t)\bigr),
 \qquad u_{q,\ell}\in\mathbb Z_\ell^\times .           \tag{E.127.6}
\]
The unit records the choice of a tame generator.
Since $\chi_\ell(h_q)=1$, the Kummer cocycle restricts
to a homomorphism on its procyclic closure.  Raising to
$e_\ell$ leaves (E.127.6) unchanged.

For our point the three valuation vectors are
\[
       v_2(t,1-t)=(-4,-4),\qquad
       v_5(t,1-t)=(1,0),\qquad
       v_{11}(t,1-t)=(0,1).                            \tag{E.127.7}
\]
If $\ell$ is odd and differs from $5,11$, use $q=5,11$.
If $\ell=5$, use $q=2,11$; their determinant is $-4$.
If $\ell=11$, use $q=2,5$; their determinant is $4$.
In every case the two vectors in (E.127.6) span
$F_\ell/\Phi(F_\ell)\simeq\mathbb F_\ell^2$.

The group $H=z_t(K)$ is closed: $K$ is compact and
$z_t|_K$ is a continuous homomorphism.  Thus
$H\cap F_\ell$ is a closed subgroup whose image in
the Frattini quotient is all of that quotient.  The
Burnside basis theorem for finitely generated pro-$\ell$
groups gives $H\cap F_\ell=F_\ell$.  Taking closures
of the products over finite sets of odd primes proves
(E.127.1).  For every $d\in D$, choose one $g$ mapping
to $d$.  Multiplication by the geometric elements just
obtained changes its odd coordinate arbitrarily while
retaining $d$.  This proves (E.127.2).

We now compute the two-primary factor, rather than only
bound it.  Let $S$ be the closed subgroup of $G_{\mathbb Q}$
generated by $\tau_5^{e_2}$ and $\tau_{11}^{e_2}$.
The group $S$ need not be pro-$2$.  Nevertheless
$\rho_2(S)=1$, so $z_{t,2}|_S$ is a homomorphism.
Its two generator images have the independent Kummer
vectors $(u_{5,2},0)$ and $(0,u_{11,2})$ modulo two.
Burnside's basis theorem therefore gives
$z_{t,2}(S)=F_2$.

The two generators of $\rho_{\rm all}(S)$ are supported
at distinct primes, respectively $5$ and $11$.  They
commute in the product of automorphism groups, so this
image is abelian.  Consequently $S'\subseteq K$, where
the prime denotes the closed commutator subgroup, and
\[
                     z_{t,2}(S')=F_2'.
                                                        \tag{E.127.8}
\]
Indeed a continuous surjection of compact groups maps
the closed commutator subgroup onto the closed commutator
subgroup: it maps the generating commutators onto all
commutators, and its compact image is closed.
The already established odd geometric factors allow us
to remove every odd coordinate from elements of $z_t(S')$.
Thus $F_2'\subseteq H_2:=H\cap F_2$.

The same elements $k_{5,2}$ and $k_{11,2}$ in
(E.127.5) have abelian coordinates
$(-2u_{5,2},0)$ and $(0,-2u_{11,2})$ by (E.127.6).
They lie in $H_2$, since each is the image under $z_t|_K$
of a pro-$2$ element.  Hence $H_2$ contains $F_2'$ and
its abelian image contains $2\mathbb Z_2^2$.  It contains
the full inverse image of $2\mathbb Z_2^2$, which is
$\Phi(F_2)$.

For the reverse inclusion, $K$ fixes every root of unity,
because each $\rho_\ell$ retains the $\ell$-adic
cyclotomic character.  The Kronecker--Weber theorem
\cite[I, Theorem~4.16, pp.49--50]{RootMilneCFT} implies
that it fixes every abelian extension of $\mathbb Q$,
hence both quadratic Kummer fields of $t$ and $1-t$.
The image of $H$ in $F_2/\Phi(F_2)$ is therefore zero.
We obtain $H_2=\Phi(F_2)$.  Since $H$ contains its entire
odd product, removing odd components from any of its
elements proves (E.127.3).

The same quadratic Kummer vanishing on $K$ defines
$\nu$ unambiguously on $D$.  It is continuous by the
compact quotient map $G_{\mathbb Q}\twoheadrightarrow D$.
The $D$-action on $F_2/\Phi(F_2)$ is trivial, so the
right side of (E.127.4) is the kernel of the surjective
homomorphism $(w,d)\mapsto\bar w_2-\nu(d)$ to
$\mathbb F_2^2$.  The actual joint image is contained
in this kernel, projects onto $D$, and has the entire
geometric kernel (E.127.3).  Choosing one lift of each
$d$ and then multiplying by that geometric kernel proves
equality.  This also proves the index-four assertion.
Changing the geometric path conjugates the image;
the Frattini subgroup and its quadratic quotient are
unchanged, so all conclusions hold for every path.
\end{proof}

\paragraph{Scope.}
This is a statement about the actual quotient images of
one rational-point section.  It supplies simultaneous
finite geometric freedom in all odd coefficient factors
while retaining all tangential arithmetic actions.  It
does not produce an ordinary infinite-support completed
pair or prove compatibility of an iterative ghost
construction.  The computed two-primary constraint is the quadratic
Kummer condition in (E.127.4), not a higher nilpotent
obstruction on a proposed ghost.
